\documentclass[11pt, a4paper]{article}
\usepackage[utf8]{inputenc}
\usepackage[T1]{fontenc}
\usepackage[spanish]{babel}
\usepackage{amsmath, amssymb}
\usepackage{geometry}
\usepackage{enumitem}
\usepackage{titlesec}
\usepackage{parskip}
\usepackage{mdframed}
\usepackage{xcolor}

\geometry{margin=2.5cm}

\titleformat{\section}{\large\bfseries}{}{0em}{}
\titleformat{\subsection}{\normalsize\bfseries}{}{0em}{}
\titleformat{\subsubsection}{\normalsize\itshape}{}{0em}{}

% Bloque de "respuesta posible"
\newmdenv[
    linecolor=gray!50,
    linewidth=0.5pt,
    backgroundcolor=gray!8,
    skipabove=6pt,
    skipbelow=6pt,
    innertopmargin=6pt,
    innerbottommargin=6pt,
]{respuesta}

\newcommand{\pregunta}[1]{\subsection{#1}}
\newcommand{\eje}[1]{\subsubsection{#1}}

\title{\textbf{Modelos de Parcial -- Neurociencia y Comportamiento} \\
\large Inteligencia Artificial y Neurociencias \\ Universidad Torcuato Di Tella}
\author{}
\date{}

\begin{document}
\maketitle

\begin{mdframed}[linecolor=black, linewidth=0.7pt]
\textbf{Formato} (igual al simulacro): en el parcial caen \textbf{cuatro preguntas}, una por cada eje tem\'atico. Responder cada una en \textbf{entre uno y tres p\'arrafos, no m\'as}. Los ejes son: (1)~Modelado computacional del comportamiento humano, (2)~Aprendizaje, (3)~Emociones, (4)~Gen\'etica del comportamiento. Cada modelo trae las preguntas y, debajo, una \emph{respuesta posible} a modo de gu\'ia.
\end{mdframed}

%% ============================================================
\section{Modelo 1}
%% ============================================================

\eje{1. Modelado computacional del comportamiento humano}
\pregunta{En el art\'iculo de Nature sobre profundidad de planificaci\'on en \textit{Four-in-a-Row}, el modelo predice tres tipos de datos observables en los humanos. \textquestiondown Cu\'ales son y qu\'e componente del modelo permite predecir cada uno?}

\begin{respuesta}
\textit{Respuesta posible.} El modelo predice tres cosas. Primero, las \textbf{elecciones} (jugadas) de las personas: validado con validaci\'on cruzada 5-fold en partidas humano vs.\ humano, predice las jugadas muy por encima del azar, gracias a la funci\'on de valor y al algoritmo de b\'usqueda que arman la variaci\'on principal. Segundo, los \textbf{tiempos de respuesta}: el tama\~no del \'arbol de decisi\'on construido por el algoritmo de b\'usqueda es predictor del tiempo de respuesta observado (\'arboles m\'as grandes, m\'as tiempo). Tercero, los \textbf{movimientos oculares}: la distribuci\'on de casillas que visita el algoritmo de b\'usqueda correlaciona con la distribuci\'on de la atenci\'on visual de los jugadores ($\rho \approx 0{,}54$). La conclusi\'on central es que los jugadores m\'as fuertes planifican m\'as profundamente y tienen menos lapsus atencionales, sin evidencia de que mejore la calidad de la heur\'istica (los pesos de las features).
\end{respuesta}

\eje{2. Aprendizaje}
\pregunta{\textquestiondown Cu\'ales son las tres condiciones que, seg\'un Darwin, deben cumplirse para que ocurra la selecci\'on natural? Explique adem\'as por qu\'e la teor\'ia de Lamarck no requiere una de ellas.}

\begin{respuesta}
\textit{Respuesta posible.} Para que ocurra selecci\'on natural deben darse tres condiciones: \textbf{variabilidad} (existen diferencias entre los rasgos de los individuos de una poblaci\'on), \textbf{heredabilidad} (esas variaciones se heredan de padres a hijos) y \textbf{adaptabilidad} (algunas de esas variaciones se adaptan mejor al ambiente, por lo que los individuos que las poseen sobreviven y se reproducen m\'as, dejando m\'as descendencia). Sin variaciones previas en la poblaci\'on, el proceso de selecci\'on no tiene sobre qu\'e operar.

En la teor\'ia de \textbf{Lamarck} los rasgos \emph{adquiridos} durante la vida del individuo se heredan a la descendencia; por eso no se necesita que existan variaciones previas en la poblaci\'on al comienzo del proceso: la variaci\'on se genera durante la vida y se transmite. En la de \textbf{Darwin}, en cambio, los rasgos adquiridos no se heredan, y las variaciones preexistentes (y heredables) son la materia prima imprescindible de la selecci\'on natural.
\end{respuesta}

\eje{3. Emociones}
\pregunta{Mencionamos cuatro afirmaciones sobre las emociones en las que la psicolog\'ia evolutiva s\'i est\'a de acuerdo. Enum\'erelas y explique brevemente la que vincula las emociones con la interocepci\'on.}

\begin{respuesta}
\textit{Respuesta posible.} Las cuatro afirmaciones son: (1)~las emociones son el resultado del \textbf{perfeccionamiento de dispositivos de supervivencia} a lo largo de la evoluci\'on de las especies; (2)~esos dispositivos usan las \textbf{mismas estructuras cerebrales y los mismos neurotransmisores en todos los mam\'iferos}, incluido el ser humano; (3)~esas estructuras est\'an vinculadas a \textbf{funciones b\'asicas de regulaci\'on del equilibrio interno}, necesario para la supervivencia; y (4)~las emociones est\'an \'intimamente arraigadas en fen\'omenos sensitivos viscerales corporales que llamamos \textbf{interocepci\'on}.

Sobre la cuarta: la interocepci\'on es la percepci\'on de los estados viscerales del organismo (pulsaciones, respiraci\'on, tensi\'on, estado del est\'omago, etc.). Las emociones no son ajenas al cuerpo, sino que est\'an enraizadas en esas se\~nales corporales. En la formulaci\'on extrema de William James, ``los cambios f\'isicos que experimentamos ante una situaci\'on \emph{son} las emociones'': la emoci\'on no antecede a la reacci\'on corporal, sino que en buena medida consiste en ella.
\end{respuesta}

\eje{4. Gen\'etica del comportamiento}
\pregunta{Explique el m\'etodo cl\'asico de comparaci\'on entre gemelos y mellizos para estimar el componente gen\'etico de un rasgo. \textquestiondown Qu\'e l\'ogica lo sustenta y cu\'al es su principal supuesto?}

\begin{respuesta}
\textit{Respuesta posible.} Los gemelos (monocig\'oticos) comparten pr\'acticamente el 100\,\% de sus alelos, mientras que los mellizos (dicig\'oticos) comparten en promedio la mitad. El m\'etodo compara la \textbf{correlaci\'on} de un rasgo entre pares de gemelos con la correlaci\'on entre pares de mellizos (criados juntos y del mismo sexo). Si la similitud entre gemelos es mayor que entre mellizos, se sospecha que el rasgo tiene un componente gen\'etico, porque la \'unica diferencia sistem\'atica entre ambos tipos de pares es cu\'antos alelos comparten. En f\'ormula:
\[
C_G = 2 \times [\mathrm{Corr}(\text{gemelos}) - \mathrm{Corr}(\text{mellizos})], \quad
C_{AC} = \mathrm{Corr}(\text{gemelos}) - C_G, \quad
C_{ANC} = 1 - C_G - C_{AC}.
\]

El \textbf{supuesto clave} es que los ambientes que generan similitudes son los mismos para gemelos que para mellizos (\emph{equal environments assumption}). Si se viola, hay sesgos: si los gemelos comparten ambientes m\'as parecidos que los mellizos (por ejemplo por tener a alguien id\'entico en el mundo), el m\'etodo \textbf{sobreestima} el componente gen\'etico; si los ambientes intrauterinos son m\'as similares entre mellizos que entre gemelos, lo \textbf{subestima}.
\end{respuesta}

%% ============================================================
\section{Modelo 2}
%% ============================================================

\eje{1. Modelado computacional del comportamiento humano}
\pregunta{Describa los tres componentes del modelo cognitivo del art\'iculo de \textit{Four-in-a-Row} y explique c\'omo es la funci\'on de valor de los estados.}

\begin{respuesta}
\textit{Respuesta posible.} Los tres componentes son una \textbf{funci\'on de valor de los estados}, un \textbf{algoritmo de b\'usqueda} y un \textbf{mecanismo de atenci\'on}. La funci\'on de valor asigna un valor a cada estado del tablero como una \textbf{suma ponderada de atributos} (features): centro, dos-en-l\'inea conectado, dos-en-l\'inea no conectado, tres-en-l\'inea y cuatro-en-l\'inea. Se computa restando el valor de las features del oponente al de las features propias, y los pesos del jugador que va a mover se multiplican por una constante de escala, para capturar que el valor de una posici\'on depende de qui\'en juegue (un tres-en-l\'inea puede ser victoria inmediata o no seg\'un a qui\'en le toque).

El \textbf{algoritmo de b\'usqueda} asume que ambos jugadores eligen la jugada que lleva al estado de mayor valor; expande el \'arbol, eval\'ua las posiciones con la funci\'on de valor y \textbf{poda} las ramas cuyo valor est\'a por debajo del mejor menos un umbral, deteni\'endose con cierta probabilidad en cada iteraci\'on (distribuci\'on geom\'etrica). El \textbf{mecanismo de atenci\'on} modela la atenci\'on selectiva descartando aleatoriamente algunas features antes de construir el \'arbol, agrega ruido gaussiano a los valores y contempla una tasa de lapsus (probabilidad de jugar al azar). El resultado: los expertos planifican m\'as profundo y tienen menos lapsus.
\end{respuesta}

\eje{2. Aprendizaje}
\pregunta{Distinga los conceptos de \textbf{rasgo adaptativo} y \textbf{subproducto}, y d\'e al menos un ejemplo de cada uno (puede ser fisiol\'ogico o comportamental). \textquestiondown Qu\'e es, adem\'as, el ``ruido'' en este esquema?}

\begin{respuesta}
\textit{Respuesta posible.} Un \textbf{rasgo adaptativo (adaptaci\'on)} es un rasgo que aument\'o el \'exito reproductivo del individuo, es decir, que fue una ventaja para la supervivencia y la reproducci\'on y por eso fue seleccionado. Un \textbf{subproducto (byproduct)} es una consecuencia de un rasgo adaptativo que no aport\'o ventaja evolutiva propia: ``viene de arrastre'' con otro rasgo que s\'i fue seleccionado.

Ejemplos fisiol\'ogicos: la \textbf{hemoglobina} (adaptaci\'on, porque permite transportar ox\'igeno) y el \textbf{color rojo de la sangre} (subproducto, es s\'olo consecuencia del color de la hemoglobina); tambi\'en el cord\'on umbilical como adaptaci\'on y el ombligo como subproducto. Ejemplos comportamentales: el \textbf{gusto por el az\'ucar}, el impulso sexual o el apego paterno-materno son adaptaciones, mientras que las adicciones, el gusto por el f\'utbol o el sexo con anticonceptivos son subproductos. El \textbf{ruido} es el tercer fruto de un proceso evolutivo: variaci\'on aleatoria que no cumple ninguna funci\'on y que simplemente aparece sin haber sido seleccionada.
\end{respuesta}

\eje{3. Emociones}
\pregunta{\textquestiondown Por qu\'e a algunas emociones se las llama \textbf{b\'asicas} (o primarias)? D\'e dos ejemplos y explique qu\'e ventaja de supervivencia o reproducci\'on se les atribuye.}

\begin{respuesta}
\textit{Respuesta posible.} Se las considera b\'asicas por tres razones: aparecen desde \textbf{edades muy tempranas} en todos los humanos (incluso en personas con ceguera y sordera cong\'enitas), son \textbf{universales} (se observan en todas las culturas, idea que ya propon\'ia Darwin) y las podemos \textbf{reconocer a partir de expresiones faciales universales}, comunes a todas las culturas e independientes del aprendizaje cultural. Ejemplos de emociones b\'asicas son el miedo, el asco, la ira, la tristeza, la alegr\'ia y la sorpresa.

Dos ejemplos con su ventaja evolutiva: el \textbf{asco} nos protege de ingerir sustancias potencialmente t\'oxicas o en mal estado (venenos, comida podrida), reduciendo el riesgo de intoxicaci\'on; el \textbf{miedo} activa el organismo para pelear o huir ante una amenaza (respuesta de lucha o huida), preparando al cuerpo para responder a situaciones peligrosas. En ambos casos son ``dispositivos de supervivencia'' que aumentaron la probabilidad de sobrevivir y dejar descendencia.
\end{respuesta}

\eje{4. Gen\'etica del comportamiento}
\pregunta{Defina el \textbf{componente gen\'etico} de un rasgo de manera rigurosa. Al estudiar el componente ambiental, \textquestiondown en qu\'e dos grandes grupos se dividen las variables y c\'omo se distingue uno de otro?}

\begin{respuesta}
\textit{Respuesta posible.} El \textbf{componente gen\'etico $(C_G)$} de un rasgo es la proporci\'on de la variaci\'on total del rasgo entre individuos que se explica por las variaciones en sus genes:
\[
C_G = \frac{V_G}{V_G + V_A} = \frac{V_G}{V_T},
\]
donde $V_G$ es la varianza gen\'etica y $V_A$ la ambiental. El componente ambiental es el porcentaje de la variaci\'on total explicado por los ambientes de desarrollo; por definici\'on, componente gen\'etico m\'as componente ambiental suman 100\,\%. Es crucial recordar que el componente gen\'etico \textbf{vale para una poblaci\'on dada en un momento dado del tiempo}: no es una propiedad fija del rasgo.

Al estudiar el componente ambiental se dividen las variables en dos grupos. Las \textbf{variables del ambiente compartido} son todo lo que no es gen\'etico y hace que los miembros de una familia se parezcan: lo que comparten dos hermanos criados juntos pero no dos separados al nacer (clase social, estilo parental, barrio de crianza). Las \textbf{variables del ambiente no compartido} son todo lo que no es gen\'etico ni se comparte entre miembros de una familia, ligado a la experiencia particular de cada individuo. En f\'ormula: $V_A = V_{AC} + V_{ANC}$.
\end{respuesta}

%% ============================================================
\section{Modelo 3}
%% ============================================================

\eje{1. Modelado computacional del comportamiento humano}
\pregunta{Seg\'un el art\'iculo, \textquestiondown en qu\'e se diferencian los expertos de los novatos? Discuta la relaci\'on entre ``planificar m\'as profundo'' y la idea de velocidad de procesamiento / reconocimiento de patrones.}

\begin{respuesta}
\textit{Respuesta posible.} El hallazgo central es que los jugadores m\'as fuertes \textbf{planifican m\'as profundamente} (exploran m\'as pasos a futuro en el \'arbol de decisi\'on) y cometen \textbf{menos lapsus atencionales} (menos jugadas al azar y menos features olvidadas). En cambio, \textbf{no} hay evidencia de que mejore la calidad de la heur\'istica: los pesos que asignan a las distintas features del tablero son similares entre expertos y novatos. Es decir, lo que distingue al experto no es ``valorar mejor'' cada posici\'on, sino buscar m\'as lejos y con menos ruido.

La literatura suele explicar la superioridad del experto por mejor \textbf{reconocimiento de patrones} y/o b\'usqueda m\'as profunda. Los autores vinculan ``planificar m\'as profundo'' con mayor \textbf{velocidad de procesamiento}: los expertos afinan su representaci\'on de las features relevantes, lo que les permite hacer m\'as evaluaciones por unidad de tiempo y por lo tanto explorar el \'arbol m\'as profundamente en el mismo lapso. Esto es consistente con un mejor reconocimiento de patrones (representaciones m\'as eficientes), pero pone el foco en un factor frecuentemente subestimado: la velocidad de procesamiento.
\end{respuesta}

\eje{2. Aprendizaje}
\pregunta{En clase vimos experimentos con beb\'es (y fetos) humanos que muestran la existencia de \emph{priors} que llevan a comportamientos instintivos. Relate dos de esos experimentos e indique qu\'e prior demuestra cada uno. \textquestiondown Para qu\'e sirven estos priors?}

\begin{respuesta}
\textit{Respuesta posible.} Un experimento ilumin\'o el vientre de embarazadas con \textbf{tres puntos de luz (l\'aseres)}. Cuando la disposici\'on se parec\'ia a un rostro ---dos puntos arriba como ojos y uno abajo como nariz---, los fetos giraban la cabeza hacia la luz con m\'as frecuencia que cuando la disposici\'on estaba invertida. Esto muestra un \textbf{prior de preferencia por las caras}, presente incluso antes de nacer. En otro tipo de experimentos, beb\'es de pocos meses miran m\'as tiempo (muestran sorpresa) ante \textbf{eventos f\'isicamente imposibles}, como una caja que parece flotar en el aire, lo que evidencia un prior que les hace inferir leyes b\'asicas de la naturaleza (los objetos no flotan, caen).

Otros priors vistos: detecci\'on r\'apida de amenazas (encontramos v\'iboras m\'as r\'apido que flores), intuici\'on num\'erica y probabil\'istica, y razonamiento l\'ogico pre-verbal. Estos comportamientos son \textbf{instintivos}: aparecen en humanos t\'ipicos con desarrollo t\'ipico, fueron seleccionados a lo largo de la evoluci\'on y sirven de \textbf{andamiaje para aprendizajes m\'as complejos} (por ejemplo, la atenci\'on innata al habla facilita el aprendizaje del lenguaje). El cerebro no es una \emph{tabula rasa}: lo innato y lo adquirido se combinan.
\end{respuesta}

\eje{3. Emociones}
\pregunta{Desde la neurociencia de la regulaci\'on emocional, explique el rol del \textbf{sistema nervioso aut\'onomo} (simp\'atico y parasimp\'atico) y en qu\'e consiste la ``receta'' para una vida larga y buena seg\'un ese enfoque.}

\begin{respuesta}
\textit{Respuesta posible.} El sistema nervioso aut\'onomo (la parte ``reptiliana'', involuntaria) tiene dos ramas. La rama \textbf{simp\'atica (SNS)} prepara al cuerpo para la acci\'on: estr\'es, atenci\'on, lucha o huida (dilata pupilas, acelera pulsaciones, libera glucosa). La rama \textbf{parasimp\'atica (PNS)} se encarga del descanso, la relajaci\'on, la digesti\'on y la recuperaci\'on, y \textbf{calma} la activaci\'on del SNS y del eje hipotal\'amico-hipofisario-adrenal (HPAA). El sufrimiento cr\'onico se asocia a una \textbf{sobreactivaci\'on del simp\'atico}, donde sensaciones corporales y pensamientos negativos se retroalimentan (los ``segundos dardos'': los pensamientos que agregamos a un dolor inevitable). En el marco de las terapias cognitivo-conductuales, Dolor $\to$ Sufrimiento est\'a mediado por el pensamiento ``terribilizador''.

La \textbf{receta} con mejores probabilidades para una vida larga y buena es mantener una activaci\'on basal predominantemente \textbf{parasimp\'atica}, con una leve activaci\'on simp\'atica para la vitalidad, y s\'olo \textbf{picos ocasionales} del simp\'atico ante grandes oportunidades o amenazas reales. La regulaci\'on emocional se entiende entonces como la t\'ecnica de ``mantener a raya'' el sistema aut\'onomo: cada vez que uno calma el ANS estimulando el parasimp\'atico, orienta cuerpo y mente hacia el bienestar y, con la repetici\'on, construye estructura neuronal (como plantea Sapolsky en \textit{\textquestiondown Por qu\'e las cebras no tienen \'ulcera?}, el estr\'es psicol\'ogico cr\'onico, exclusivo de los humanos, genera enfermedad).
\end{respuesta}

\eje{4. Gen\'etica del comportamiento}
\pregunta{\textquestiondown Qu\'e es una \textbf{correlaci\'on genotipo-ambiente}? Describa sus tres tipos (pasiva, evocativa y activa) con un ejemplo de cada uno, y explique por qu\'e esto complica la interpretaci\'on causal de estudios que correlacionan ambientes con rasgos.}

\begin{respuesta}
\textit{Respuesta posible.} Una \textbf{correlaci\'on genotipo-ambiente} ocurre porque los \textbf{genes influyen sobre los ambientes} a los que un individuo queda expuesto: no s\'olo genes y ambientes influyen sobre los rasgos, sino que los genes moldean el propio ambiente. Hay tres tipos. \textbf{Pasiva}: el ni\~no recibe de sus padres genes que correlacionan con el ambiente que esos mismos padres crean (p.\ ej., si unos alelos predisponen a la lectura, es probable que los padres tengan esos alelos y por eso haya m\'as libros en la casa). \textbf{Evocativa}: el ni\~no es tratado seg\'un sus propensiones gen\'eticas (un chico con facilidad para los n\'umeros recibe m\'as regalos con n\'umeros). \textbf{Activa}: el ni\~no busca y crea activamente ambientes acordes a sus propensiones (ese mismo chico elige juguetes con n\'umeros).

Esto complica la interpretaci\'on causal porque una correlaci\'on observada entre un ambiente y un rasgo puede estar \textbf{mediada por la gen\'etica} y no ser causal. Ejemplo del curso: se observa que en familias donde los padres incentivan la curiosidad los hijos son m\'as curiosos; pero esto \textbf{no es evidencia causal}, porque puede tratarse de una correlaci\'on pasiva (los hijos heredan las propensiones de los padres y a la vez los padres moldean el ambiente). Un caso real: la asociaci\'on entre abandono paterno / tabaquismo materno en el embarazo y depresi\'on de los hijos desaparece al controlar por el puntaje polig\'enico de los padres, lo que sugiere una correlaci\'on genotipo-ambiente pasiva y no una relaci\'on causal directa.
\end{respuesta}

%% ============================================================
\section{Banco de preguntas extra (para practicar)}
%% ============================================================

\subsection*{Modelado computacional}
\begin{itemize}[nosep]
    \item \textquestiondown Por qu\'e \textit{Four-in-a-Row} es un buen juego para medir profundidad de planificaci\'on y el ajedrez no? (tratabilidad computacional, reglas simples, novedad, complejidad suficiente para distinguir expertos de novatos).
    \item \textquestiondown Qu\'e es la ``variaci\'on principal'' en el \'arbol de decisi\'on y c\'omo funciona la poda?
    \item Hip\'otesis de implementaci\'on neural del modelo: \textquestiondown qu\'e regi\'on se asocia a cada componente? (funci\'on de valor $\to$ corteza orbitofrontal; b\'usqueda $\to$ prefrontal, c\'ingulo, hipocampo; atenci\'on $\to$ red frontoparietal).
\end{itemize}

\subsection*{Aprendizaje}
\begin{itemize}[nosep]
    \item \textquestiondown Qu\'e es la \textbf{teor\'ia de la carga cognitiva} y qu\'e nos hace expertos seg\'un ella? (la memoria de trabajo es un cuello de botella por su baja capacidad y durabilidad; la memoria de largo plazo, de capacidad casi ilimitada, es la que nos hace pensar mejor y no puede sustituirse por el acceso r\'apido a informaci\'on externa como Google o una IA).
    \item Las tres etapas de todo aprendizaje: adquisici\'on, fluidez y generalizaci\'on (m\'as mantenimiento). Describa cada una.
    \item Los cuatro pilares del aprendizaje de Dehaene: atenci\'on, feedback (\textexclamdown y regresi\'on a la media!), motivaci\'on (intr\'inseca vs.\ extr\'inseca) y consolidaci\'on.
    \item Explique la \textbf{regresi\'on a la media} y por qu\'e genera la ilusi\'on de que castigar mejora y premiar empeora.
    \item \textquestiondown Qu\'e relaci\'on hay entre la estructura del cerebro (hiperpar\'ametros) seleccionada por evoluci\'on y el \emph{fine tuning} durante la vida?
\end{itemize}

\subsection*{Emociones}
\begin{itemize}[nosep]
    \item Diferencia entre \textbf{emoci\'on} y \textbf{sentimiento} (los sentimientos se asocian m\'as a lo cognitivo y duran m\'as en el tiempo).
    \item \textquestiondown Qu\'e son los ``segundos dardos'' y c\'omo se relacionan Dolor y Sufrimiento?
    \item \textquestiondown Por qu\'e ``las cebras no tienen \'ulcera''? (el estr\'es animal es agudo y puntual; el humano inventa estresores psicol\'ogicos cr\'onicos que enferman).
\end{itemize}

\subsection*{Gen\'etica del comportamiento}
\begin{itemize}[nosep]
    \item Diferencia entre las dos definiciones de ``rasgo gen\'etico'': la de la psicolog\'ia evolutiva (lo que todos tenemos en com\'un, p.\ ej.\ aprender un lenguaje) y la de la gen\'etica del comportamiento (lo que nos hace diferentes, las variaciones individuales).
    \item M\'etodo cl\'asico con \textbf{hijos adoptivos}: comparar correlaci\'on padres-hijos biol\'ogicos vs.\ padres-hijos adoptivos. \textquestiondown Qu\'e estima cada comparaci\'on? Problemas: colocaci\'on selectiva, ambiente intrauterino, generalizaci\'on.
    \item \textquestiondown Qu\'e es un \textbf{puntaje polig\'enico} y qu\'e es la \textbf{heredabilidad faltante}?
    \item \textquestiondown Qu\'e es la \textbf{interacci\'on genotipo-ambiente} (distinta de la correlaci\'on)? Ejemplo del estudio de maltrato infantil y TDAH: las rectas paralelas indican que \textbf{no} hay interacci\'on (el maltrato aumenta s\'intomas independientemente del puntaje polig\'enico).
    \item Componente gen\'etico de rasgos concretos: autismo $\sim$80\,\%, esquizofrenia $\sim$50\,\%, depresi\'on $\sim$37--48\,\%, IMC $\sim$40--80\,\%, zurdera $\sim$25\,\%, extroversi\'on $\sim$37\,\%.
\end{itemize}

\end{document}
